(1) 三维声学频率空间中近似态密度 $g(\omega)=A\omega^2$，计数要求 $3N=\int_0^{\omega_D}g(\omega)d\omega=A\omega_D^3/3$，故 $A=9N/\omega_D^3$。德拜频率通过总模数保证正确的自由度数，而非任选的截断常数。
(2) 声子为玻色准粒子，平均占据数 $n_B(\omega)=[e^{\hbar\omega/(k_BT)}-1]^{-1}$，因此热激发能 $U(T)=\int_0^{\omega_D}(9N/\omega_D^3)\hbar\omega^3n_B(\omega)d\omega$。置 $x=\hbar\omega/(k_BT)$、$\Theta_D=\hbar\omega_D/k_B$，得 $U=9Nk_BT(T/\Theta_D)^3\int_0^{\Theta_D/T}x^3/(e^x-1)dx$。
(3) 在 $T\ll\Theta_D$ 时积分上界可推至无穷，且 $\int_0^\infty x^3/(e^x-1)dx=\pi^4/15$，因此 $U\simeq(3\pi^4/5)Nk_BT^4/\Theta_D^3$。定容求导得 $C_V\simeq(12\pi^4/5)Nk_B(T/\Theta_D)^3$，这就是三维声学模导致的德拜 $T^3$ 定律。
(4) 高温 $\hbar\omega\ll k_BT$，每个模式平均热能约 $k_BT$，从而 $U\simeq3Nk_BT$、$C_V\to3Nk_B$，即杜隆–珀蒂定律。低温极限依赖线性色散与态密度 $\omega^2$；高温极限仅适用于远低于材料熔化、非谐效应不严重的晶格热容，不可把声子简并当作恒定的三自由度粒子。
